% Figure: a block pushed by a horizontal force near its top, illustrating
% the competition between sliding (friction limit) and tipping (moment
% about the leading bottom edge). Two reactions are drawn: the friction at
% the base and the leading-edge pivot the block tips about.
\begin{figure}[htbp]
\centering
\begin{tikzpicture}[
  push/.style={draw=engred, -{Stealth}, very thick},
  wt/.style={draw=engblue, -{Stealth}, very thick},
  nrm/.style={draw=engamber, -{Stealth}, very thick},
  fric/.style={draw=engorange, -{Stealth}, very thick},
  dim/.style={draw=enggray, thin, {Stealth}-{Stealth}},
  scale=1.0,
]
% ground
\draw[draw=engblack, very thick] (-1.0,0) -- (5.6,0);
\foreach \x in {-0.7,-0.2,...,5.3} {\draw[draw=enggray, thin] (\x,0) -- (\x-0.2,-0.25);}
% the block: width b, height H
\draw[draw=engblack, fill=engblue!8, thick] (1.2,0) rectangle (2.6,2.6);
% centre of mass
\coordinate (G) at (1.9,1.3);
\fill[engblack] (G) circle (0.06);
\node[font=\scriptsize, anchor=south east] at (G) {$G$};
% weight at G
\draw[wt] (G) -- ++(0,-1.0) node[anchor=north east, font=\scriptsize, engblue] {$mg$};
% applied push near the top, at height h
\coordinate (Ppt) at (1.2,2.2);
\draw[push] ($(Ppt)+(-1.1,0)$) -- (Ppt) node[midway, anchor=south, font=\scriptsize, engred] {$P$};
\node[font=\scriptsize, anchor=east, engred] at ($(Ppt)+(-1.1,0)$) {};
% normal reaction (resultant) shifted toward the leading edge
\draw[nrm] (2.3,0) -- ++(0,1.3) node[anchor=south, font=\scriptsize, engamber] {$N$};
% friction at the base, opposing the push
\draw[fric] (1.9,0.0) -- ++(0.9,0) node[anchor=west, font=\scriptsize, engorange] {$F_f$};
% leading bottom edge: the tipping pivot
\fill[engred] (2.6,0) circle (0.07);
\node[font=\scriptsize, anchor=north west, engred] at (2.6,0) {$O$};
% dimension hints: width b and push height h
\draw[dim] (1.2,-0.55) -- (2.6,-0.55);
\node[font=\scriptsize, anchor=north] at (1.9,-0.55) {$b$};
\draw[dim] (3.0,0) -- (3.0,2.2);
\node[font=\scriptsize, anchor=west] at (3.0,1.1) {$h$};
% a faint tipping arc about O
\draw[draw=engred, thin, dashed] (1.2,2.6) arc (143:175:3.2);
\node[font=\scriptsize, anchor=south west, engred] at (4.1,1.7) {tips about $O$};
\end{tikzpicture}
\caption[Tipping versus sliding of a pushed block]{A block of width $b$
and weight $mg$ pushed by a horizontal force $P$ applied at height $h$.
The chosen body is the block. The base supplies a distributed reaction
whose resultant is a normal force $N$ (which migrates toward the leading
edge as $P$ grows) and a friction force $F_f$ opposing the push. Two
failure modes compete. Sliding begins when $P$ reaches the friction
limit, $P = \mu_s mg$, independent of where $P$ is applied. Tipping
begins when $P$ generates enough moment about the leading bottom edge
$O$ to lift the trailing edge, $P h = mg\,(b/2)$, which depends on the
push height $h$. Whichever threshold is lower governs: a high push on a
high-friction floor tips, a low push on a slick floor slides. The two
lines cross at $h = b/(2\mu_s)$, the geometry at which the block is
equally close to both failures.}
\label{fig:vol03ch01:tipping-vs-sliding}
\end{figure}
